\documentclass [11pt,a4paper,oneside,openany]{book} %Classe del documento, formato carta, singola facciata, apertura capitoli pag destra/sinista (ininfluente se impostato oneside). Formato libro
%%\usepackage[italian]{babel} %Lingua Documento (da impostare per la sillabazione)
\usepackage{graphicx} %Pacchetto necessario per la gestione delle immagini
\usepackage{xcolor,colortbl} %Colori
\usepackage {fancyhdr} %Pacchetto per la gestione accurata della pagina
\usepackage {setspace} %Pacchetto necessario per i comandi successivi onehalfspacing, singlespacing....
\usepackage{titlesec} % Pacchetto per le sezioni e sottosezioni
\usepackage{longtable} %Pacchetto per la gestione delle tabelle grandi
\usepackage[colorlinks=true]{hyperref} %Segnalibri nel pdf finale, sotto relativi parametri
\usepackage[T1]{fontenc}
\usepackage[utf8]{inputenc} %Pacchetti per le lettere accentate utf8
\usepackage{listings} %Pacchetto per formattare codice sorgente
\usepackage{mathtools} %Pacchetto matematica
\usepackage{float} %Pacchetto per il posizionamento delle figure

\hypersetup{
	bookmarksnumbered=true,
	linkcolor=black,
	citecolor=black,
	urlcolor=black,
}

\usepackage{geometry} % Dimensione pagina

\geometry{a4paper} % Formato carta
\addtolength{\textheight}{75pt} %Margini
\oddsidemargin 30pt

\setcounter{secnumdepth}{2}
\setcounter{tocdepth}{3}


\newcommand{\reqref}[1]{\textbf{R\ref{#1}}}
\newcommand{\goalref}[1]{\textbf{G\ref{#1}}}


\begin{document} %Inizio Documento

%%FRONTESPIZIO%%=========================================
\begin{titlepage}

 \begin{center} 
     \includegraphics[height=6cm]{logo.eps}\\
     \vspace{4em}
     {\Large \textsc{Software Engineering II}}\\
     \vspace{6em}
     {\LARGE \textbf{Travlendar+}}\\
     \vspace{3em}
     {\Large \textsc{Requirements Analysis}}\\
     \vspace{1em}
     {\Large \textsc{and}}\\
     \vspace{1em}
     {\Large \textsc{Specifications Document}}\\
 \end{center}
 
    \vskip 2cm
 
 	Authors:
 	\vspace{0.5em}
 	\begin{center}
      {\Large \textit{Edoardo D'Amico}}\\
      {\Large \textit{Gabbolini Giovanni}}\\
      {\Large \textit{Parroni Federico}}\\
    \end{center}

\vskip 2.0cm
\begin{center}
{\normalsize 1$^{st}$ October 2017}
\end{center}

\end{titlepage}

\newpage

%%INTESTAZIONI PAGINA%%==================================
\pagestyle{fancy}
\renewcommand{\chaptername}{Chapter}
\renewcommand{\chaptermark} [1]{\chaptername\ \thechapter.\ #1}{} 
\renewcommand{\chaptermark}[1]{\markboth{\thechapter.\ #1}{}} 
\renewcommand{\sectionmark}[1]{\markright{\thesection\ #1}}
\fancyhf{}
\fancyhead[LE,RO]{\bfseries\thepage} 
\fancyhead[LO,RE]{\bfseries\leftmark} 
\fancypagestyle{plain}{%
\fancyhead{} % leva l'intestazione
\renewcommand{\headrulewidth}{0pt} % e la linea
}


%%INDICE%%==============================================
\renewcommand{\contentsname}{Indice}
\renewcommand{\listfigurename}{Elenco delle Figure}
\renewcommand{\listtablename}{Elenco delle Tabelle}
\tableofcontents

\linespread{1.5}\selectfont		%interlinea
%%CAPITOLI===============================================
%%CAPITOLO 1: Introduzione=======================================
\chapter{Introduction}

\section{Purpose}
The purpose of this document is to provide more technical details than the RASD about Travelendar+ system. This document is addressed to developers and aims to identify:
\begin{itemize}
\item The high level architecture
\item The design patterns
\item The main components and their interfaces provided one for another
\item The Runtime behavior
\end{itemize}

\section{Scope}

Travlendar+ is an application that aims to ease the management of the daily appointments of a registered user by providing a desired schedule for those. 
The application runs within the context of a restricted geographic area and thus should consider only the travel means available in this scope.
The application is able to interact with different data sources, in particular Travel Means APIs and User device GPS, and to arrange the appointments by exploiting these information.
Once the schedule is computed, the application gives to the user the possibility to buy tickets for the involved travel means, then he/she can run it and follow directions provided by the application. 

\section{Definitions, Acronyms, Synonims}

\subsection{Definitions}
\theoremstyle{definition}
\newtheorem{definition}{Definition}[section]
 
\begin{definition} \label{def:userdata}
The user data include:
\begin{itemize}
\item Schedules;
\item Appointments;
\item User parameters.
\end{itemize}

\end{definition}
  
\begin{definition} \label{def:device}
A device is a PC, a Tablet or a Smartphone in which run the last version of his O.S.;
\end{definition}

\begin{definition} \label{def:schedule}
A Schedule is a set of time-ordered and not overlapping appointments where their starting times are fixed and they're linked to each other by a path covered with a specific transportation mean;
\end{definition}

\begin{definition} \label{def:bestTravelMean}
The bestTravelMean is the travel mean that satisfies the constraints on the appointment and on the schedule and that optimize the chosen criteria.
\end{definition}


\subsection{Acronyms}
%Acronyms used in the text:
\begin{itemize}
\item API: Application Programming Interface;
\item CRUD: Create/Read/Update/Delete;
\end{itemize}

%\reqref{req:R1}
%\goalref{goal:G1}

\section{Revision history}

\section{Reference Documents}
RASD document

\section{Document structure}
\begin{itemize}
\item Introduction: this section introduces the design document. It contains a justification of his utility
\item Architecture Design: this section is divided into different parts:
\begin{enumerate}
\item Overview: this section explains the division in tiers of our application and lists further details that were not specified in the RASD.
\item Components View: this sections gives a global view of the components of the application, other than a brief description, and describes how they communicate
\item Deployment view: this section shows the components that must be deployed to have the application running correctly
\item Runtime view: sequence diagrams are represented in this section to show how the components interact and collaborate during the fullfilling of a specific task. Each task is associated to each goal specified on the RASD document
\item Component interfaces: the interfaces between the components are presented in this section
\item Selected architectural styles and patterns : this section explains the architectural and design choices taken during the creation of the application
\end{enumerate}
\item Algorithms Design: this section describes the most important algorithmic parts. Pseudo code is used in order to hide unnecessary implementation details and abstract from a specific programming language
\item Requirements Traceability: this section aims to explain how the decisions taken in the RASD are linked to design elements
\end{itemize}



 
 %nome del file LaTeX del capitolo
\chapter{Overall Description}

\section{Product Perspective}

\subsection{User Model}
A user is represented within the application by a set of parameters:
Some important informations about the user are held by the following ones:
\begin{itemize}
\item \textit{travelPass}: indicates if the user has a pass for public transportation;
\item \textit{hasBike};
\item \textit{hasCar}.
\end{itemize}

\subsection{Appointment Model}
An appointment is represented within the application by a set of parameters:
\begin{itemize}
\item \textit{duration}: the time extension of the appointment
\item \textit{startingTime} or \textit{timeInterval}: the first should be given if the starting hour is well-known (deterministic), otherwise a time interval in which the appointment will be held it's provided. 
%In particular: \begin{equation}timeInterval.start + duration \leq timeInterval.end \label{eqn:constrOnTimeInterval} \end{equation}
\item \textit{location}: identifies the coordinates of the place where the appointment will be held;
\item \textit{recurrent}: specifies if the appointment will be repeated over a fixed period of time;
\item \textit{peopleVariation}: represents a variation occourring when the user picks up or leaves off someone.
\end{itemize}

The life cycle of an appointment can be represented by the following statechart:
\begin{figure}[H]
\begin{center}
\includegraphics[width=\textwidth, keepaspectratio]{"images/stateChartAppointment".png}
\caption{Appointment statechart}
\end{center}
\end{figure}

A newly-created appointment is \textbf{Unscheduled}. It could remain \textbf{Unscheduled} either when edited or there isn't a possible arrangement when a schedule is performed. Otherwise it becomes \textbf{Scheduled} if there's a feasible way to arrange it. When a scheduled appointment is edited all the appointments in that schedule return to be \textbf{Unscheduled}, because it's possible that they bring to a different schedule. When a scheduled appointment is fullfilled it becomes \textbf{Archived} and stored in the schedule history. If this last one is a recurrent appointment it must be reinsert in the list of unscheduled appointment so it will become \textbf{Unscheduled} again. The user can cancel an appointment in every moment. 

\subsection{Schedule Model} \label{subsect:schmodel}
A schedule is a set of Appointments in a given day, ordered by the scheduler following the criteria described below.
A schedule is characterized by the following variables:
\begin{itemize}
\item \textit{date};
\item \textit{startingPosition}: is the starting location of the user (e.g. user's home);
\item \textit{startingNumberOfPeople}: the number of people that must reach the first appointment.
\item \textit{wakeUpTime}: it is the starting time from which the schedule should start arranging appointments.
\end{itemize}

\subsubsection{The optimization criteria} \label{subsubsect:optcriteria} 
The criteria that can be chosen for a schedule by the user  for the optimization are the following:
\begin{itemize}
\item \textit{Minimize carbon footprint}: the scheduler will try to minimize the amount of kilometers travelled in polluting means;
\item \textit{Minimize money spent}: the scheduler will try to avoid expensive means and to exploit the public ones (especially if the user has a pass) or going by bike or on foot;
\item \textit{Minimize travelling time}: the scheduler will compute the quickest possible path reaching all the appointments locations.
\end{itemize}

\subsection{Constraints}
Constraints are impositions on some parameters managed by the system during the process of scheduling the appointments. We can distinguish between constraints on schedule and contraints on the single appointment. These can be selected by the user when he inserts an appointment or when he requests a schedule, otherwise the constraints are initialized to default values. 

\subsubsection{Constraints on schedule} 
\begin{itemize}
\item \textit{Maximum travelling distance with a specific travel mean}: the user can set a maximum amount of km to travel with a travel mean;
\item\textit{ Travel means time slots}: user can specify a time interval in which a travel mean can be used;
\item \textit{User can deactivate a particular travel means};
\item \textit{User can select which travel means uses under certain weather condition}.
\end{itemize}

\subsubsection{Constraints on appointment}
\begin{itemize}
\item \textit{User can deactivate a particular travel means}.
\end{itemize}

\subsection{Class Diagram}
\begin{figure}[H]
\begin{center}
\includegraphics[width=\textwidth, keepaspectratio]{"images/classDiagram".png}
\caption{System Class Diagram}
\end{center}
\end{figure}

\section{Product Functions}
The following requirements are derived in order to fullfill the specified goals.
\\

Requirements for \goalref{goal:G1}:
\begin{enumerate}
\renewcommand\labelenumi{\textbf{R\theenumi}}
\item The system SPW to create a new appointment allowing user to specify all its parameters; \label{req:R1}
\item The system SPW to add constraints to a previously created appointment specifying which travel means have to be avoided during the travel to the meeting location; \label{req:R2}
\\

Requirements for \goalref{goal:G2}:

\item Allow the user to set the parameters of the schedule (\ref{subsect:schmodel}), or to accept the default values;
 
\item Allow the user to select the optimization criteria (\ref{subsubsect:optcriteria}) for the schedule;

\item The system S.B.A. to gather information from external APIs about: 
\begin{itemize}
\item travel options with related travel option data;
\item weather forecast;
\item strike days;
\end{itemize}

\item The system S.B.A. to select the best travel option according to the optimization criteria taking into account:
user constraint
travel option data 
weather forecast
information about strike day


%\item The system S.B.A. to compare different travel options with related travel option data according to the optimization criteria;

\item The system S.B.A. to avoid a particular travel mean under a certain weather condition;

\end{enumerate}

\subsection{ Assumptions, dependencies and constraints}

\begin{itemize}
\item The system should be able to retrieve information about public travel means. In particular:
\begin{itemize}
\item information about delays;
\item information on possible strike days.
\end{itemize}

\item The system should be able to retrieve information about shared means. In particular:
\begin{itemize}
\item position of the available ones;  
\item prices per time unit;
\end{itemize}

\item It's possible to retrieve all travel options and travel option data from external APIs;
\end{itemize}

%list of requirements given the goals



%input{cap3_struttura_del_progetto}
%\input{cap4_calibrazione_della_videocamera}
%\input{cap5_riconoscimento}
%\input{cap6_localizzazione}
%\input{cap7_inseguitore}
%\input{cap8_test}
%\input{cap9_conclusioni}
%...


%%APPENDICI==============================================
\appendix
%%\input{appendiceA}

\end{document}